\section{Executive Summary}

\subsection{Goal and Current Position}

This branch proposes a profile-driven direction for \UML{} while
preserving the existing contract: the upstream Linux kernel built as
\texttt{arch/um} and executed as a host process. The target is not
another virtual machine monitor, not a container runtime, and not a
Linux reimplementation. The branch-local work is the combination of
swappable execution backends, runtime instrumentation gates,
profile-driven builds, and validation artifacts from one source tree.

\begin{figure}[h]
\centering
\begin{tikzpicture}[
  node distance=0.55cm,
  box/.style={rounded corners=2pt,draw=UMLGrid,fill=UMLPanel,text width=0.84\textwidth,align=left,inner sep=8pt},
  title/.style={font=\bfseries\large,text=UMLInk},
  key/.style={font=\bfseries,text=UMLBlue},
]
\node[box] (vision) {\textbf{Branch claim.} One kernel tree can produce fast production \UML{}, sanitizer-heavy research \UML{}, fuzzing \UML{}, sandbox \UML{}, library-mode \UML{}, and time-travel \UML{}.};
\node[box,below=of vision] (arch) {\textbf{Architecture.} Layer 1 abstracts the trap mechanism with \texttt{struct um\_backend\_ops}; Layer 2 places static-key gates on hot paths; Layer 3 composes compile-time sanitizers and profile features.};
\node[box,below=of arch] (state) {\textbf{Current state at \ReportDate.} Backend abstraction, static-key infrastructure, profile work, \texttt{umlctl}, selftests, and \KVMvTwo{} are in-tree. \KVMvTwo{} has working guest entry, syscall dispatch, exceptions, signals, SMP state repair, and post-gadget performance; sustained validation is in progress.};
\node[box,below=of state] (evidence) {\textbf{Evidence.} \KVMvTwo{} matches seccomp on the substrate gate, reaches cpython-parity 21/21 under UP and SMP, passes 200/200 \KVMvTwo{} iterations in the first post-T57 two-hour soak, and is faster than seccomp on all eight benchmark hosts.};
\end{tikzpicture}
\caption{The report in four claims. Each claim maps back to a tracked source document or measured artefact.}
\end{figure}

\subsection{Existing UML Contract and Branch Questions}

The branch assumes the existing \UML{} contract: Linux itself runs as a
userspace process, remains debuggable and rootless, and can host
isolated Linux instances without emulating a full machine
\cite{dike-als2000,uml-howto,uml-homepage}. The proposed changes keep
that contract and concentrate on branch-local pressure points:
userspace trap cost, multiprocessor state ownership, profile
composition, and validation breadth.

\begin{table}[h]
\centering
\caption{Existing UML contract mapped to the branch-local response.}
\small
\begin{tabularx}{\linewidth}{@{}L{0.24\linewidth}Y Y@{}}
\toprule
Existing property & Constraint carried forward & Branch-local response \\
\midrule
Kernel as a normal host process & Preserve rootless development, CI usability, and reproducible debugging where nested virtualization is unavailable. & Keep stock-host execution through ptrace/seccomp, but add KVM as a fast backend when available. \\
Isolated Linux instances & Preserve cheap Linux environments without new hardware or privileged host setup. & Make sandbox a profile, not the whole architecture; choose minimal code, seccomp-only defaults, and explicit validation. \\
Faithful upstream kernel & Preserve Linux behavior rather than introducing a userspace reimplementation with semantic drift. & Build the upstream kernel as \UML{}; profiles change defaults and instrumentation, not kernel semantics. \\
Debuggable experimentation & Preserve normal \UML{} debugging value while adding tracing, coverage, and replay hooks as explicit features. & Layer 2 static keys make observability a runtime decision; Layer 3 composes sanitizers by profile. \\
Known hard problems & Treat trap cost and SMP state ownership as implementation and validation problems, not as presentation claims. & \KVMvTwo{} and sustained validation directly target syscall cost, SMP state ownership, parity, and long-run reliability. \\
\bottomrule
\end{tabularx}
\end{table}

\begin{table}[h]
\centering
\caption{Questions the report uses to bound the branch claims.}
\small
\begin{tabularx}{\linewidth}{@{}L{0.04\linewidth}Y Y@{}}
\toprule
\# & Question & Where answered \\
\midrule
1 & What branch-local value proposition is being proposed? & Branch claim, non-goals, existing-contract table, and profile matrix. \\
2 & What is the boundary relative to QEMU, Firecracker, gVisor, LKL, and containers? & Non-goals, design boundary, and adjacent-work comparisons. \\
3 & What single architecture lets fast production, research, fuzzing, sandbox, library, and time-travel coexist? & Three-layer architecture and profile sections. \\
4 & What has actually landed, and what is still only planned? & Current-state tables, source map, and upstream queue. \\
5 & Does \KVMvTwo{} preserve Linux-visible state across SMP, signals, FPU/XSAVE, faults, and task switches? & KVM design, state-audit method, recent-fix table, and validation ladder. \\
6 & Does the performance win survive beyond a microbenchmark? & Micro, Python, and stress-tier benchmark sections. \\
7 & What validation bar remains before upstreaming? & Sustained validation, Tier 2/3/LTP, and next-work sections. \\
8 & What should maintainer review put pressure on next? & Open-state list, remaining risks, and upstream-path sections. \\
\bottomrule
\end{tabularx}
\end{table}

\subsection{Cutoff Snapshot}

\begin{figure}[h]
\centering
\metriccard{UMLBlue}{21/21}{cpython parity\\UP and SMP}\hfill
\metriccard{UMLGreen}{200/200}{\KVMvTwo{} two-hour\\validation pilot}\hfill
\metriccard{UMLAmber}{193--908x}{micro speedup\\over seccomp}\hfill
\metriccard{UMLTeal}{7 series}{staged upstream\\submission queue}
\caption{Dashboard metrics at the reporting cutoff. The numbers are intentionally sparse: each is backed by a source table later in the report.}
\end{figure}

\begin{table}[h]
\centering
\caption{Reporting snapshot.}
\begin{tabularx}{\linewidth}{@{}>{\bfseries}L{0.27\linewidth}Y@{}}
\toprule
Reporting date & \ReportDate \\
Cutoff commit & \texttt{\CutoffCommit} on \texttt{\BranchName} \\
Project corpus & 253 files under \texttt{Documentation/virt/uml/redesign/}; 32 top-level \texttt{tools/testing/selftests/um/} test areas; \texttt{arch/um/backend/kvm-v2/} implementation present \\
Functional headline & \KVMvTwo{} has working guest entry, syscall dispatch, exceptions, signals, SMP state repair, and post-gadget performance; formal 24-hour validation remains open \\
Parity headline & cpython-parity is 21/21 under both UP and SMP; substrate gate is PASS=25/FAIL=3/XFAIL=3, matching seccomp \\
Performance headline & post-gadget \KVMvTwo{} beats seccomp on all eight lab hosts: micro 193--908x, Python 2.99--4.06x, stress 3.71--8.72x \\
Validation headline & first post-T57 two-hour soak: \KVMvTwo{} 200/200 PASS; aggregate 397/400 PASS, with the three failures isolated to iocheck/seccomp timeout behavior \\
\bottomrule
\end{tabularx}
\end{table}

\subsection{What Is Done, What Is Not}

\begin{table}[h]
\centering
\caption{Status by major workstream.}
\begin{tabularx}{\linewidth}{@{}p{0.10\linewidth}p{0.22\linewidth}p{0.18\linewidth}X@{}}
\toprule
ID & Workstream & State & Summary \\
\midrule
A & Backend abstraction & \donepill & \texttt{um\_backend\_ops} contract, ptrace/seccomp split, KUnit contract, and documentation landed. \\
B & Static-key hot paths & \donepill & Hot-path gates, debugfs controls, and runtime flip demonstration landed. \\
C & Profiles and gaps & \statuspill{UMLGreen}{MOSTLY} & Defconfigs and major sanitizer/tracing ports are largely landed; snapshot, syzkaller, and some resumed features remain. \\
D & KVM backend & \progresspill & \KVMvTwo{} is functional and performant; 24-hour validation, Tier 2/3, LTP, and upstream packaging are the active edge. \\
\bottomrule
\end{tabularx}
\end{table}

\subsection{First Three-Page Takeaway}

The branch is best understood as a transition from one default
execution story to an explicit profile matrix. The common lower layer
is a backend contract; production can choose \KVM{}, CI can choose
\Seccomp{}, legacy hosts can choose ptrace, and each profile can decide
whether it wants no hooks, runtime hooks, or heavy compile-time
instrumentation.

The current branch has crossed the highest-risk implementation path:
\KVMvTwo{} can run real guest userspace, handle SMP, preserve
cross-task state, and beat \Seccomp{} on measured workloads. The
remaining gate is sustained validation: reliability, breadth, and
upstreamable series boundaries.
